\section{Distributions}\label{chap:distributions}\label{chap:fourier-transform}

\subsection{Test Functions and Distributions}\label{sec:test-functions-distributions}

\subsubsection{Motivation}\label{subsec:distributions-motivation}

Many problems arising naturally in differential equations call for a generalized definition of functions, derivatives, convergence, integrals, etc. In this subsection, we discuss a number of such questions, which will be adequately answered below.

\begin{enumerate}
\item In Chapter~\ref{chap:introduction}, we noted that any twice differentiable function of the form $u(x,t)=F(x+t)+G(x-t)$ is a solution of the wave equation $u_{tt}=u_{xx}$. Clearly, it seems natural to call $u$ a "generalized" solution even if $F$ and $G$ are not twice differentiable. A natural question is what meaning can be given to $u_{tt}$ and $u_{xx}$ in this case; obviously, they cannot be "functions" in the usual sense. The same question arises for the shock solutions of hyperbolic conservation laws which we discussed in Chapter~\ref{chap:conservation-laws-shocks}.
\item Consider the ODE initial-value problem
\begin{equation}\label{eq:impulse-ode-initial-value-problem}
u'(t)=f_\epsilon(t),\ u(0)=0,
\end{equation}
where
\begin{equation}\label{eq:unit-impulse-approximating-pulse}
f_\epsilon(t)=\begin{cases}\dfrac1\epsilon&1<t<1+\epsilon\\0&\text{otherwise.}\end{cases}
\end{equation}
Obviously, the solution is
\begin{equation}\label{eq:impulse-ode-piecewise-solution}
u(t)=\begin{cases}0&0\leq t\leq1\\\dfrac{t-1}{\epsilon}&1\leq t\leq1+\epsilon\\1&t\geq1+\epsilon.\end{cases}
\end{equation}
Note that the limit of $u$ as $\epsilon\to0$ exists; it is a step function. The function $f_\epsilon$ has unit integral; it is supported on shorter and shorter time intervals as $\epsilon$ tends to zero. It would be natural to regard the "limit" of $f_\epsilon$ as an instantaneous unit impulse. The question arises what meaning can be given to this limit and in what sense the differential equation holds in the limit. Similar questions arise in many physical problems involving idealized point singularities: the electric field of a point charge, light emitted by a point source, etc.
\item In Chapter~\ref{chap:introduction}, we outlined the solution of Dirichlet\textquotesingle s problem by minimizing the integral $\int_\Omega|\nabla u|^2\,\mathrm{d}\bm{x}$. A fundamental ingredient in turning these ideas into a rigorous theory is obviously the definition of a class of functions for which the integral is finite; the square root of the integral naturally defines a norm on this space of functions. It turns out that $C^1(\overline\Omega)$ is too restrictive; it is not a complete metric space in the norm defined by the integral. It is natural to consider the completion; this leads to functions for which $\nabla u$ does not exist in the sense of the classical definition as a pointwise limit of difference quotients.
\item The Fourier transform is a natural tool for dealing with PDEs with constant coefficients posed on all of space. However, the class of functions for which the Fourier integral exists in the conventional sense is rather restrictive; in particular, such functions must be integrable at infinity. Clearly, it would be useful to have a notion of the Fourier transform for functions which do not satisfy such a restriction, e.g., constant functions.
\end{enumerate}

The idea behind generalized functions is roughly this: Given a continuous function $f(\bm{x})$ on $\Omega$, we can define a linear mapping
\begin{equation}\label{eq:function-induced-test-functional}
\phi\mapsto\int_\Omega f(\bm{x})\phi(\bm{x})\,\mathrm{d}\bm{x}
\end{equation}
from a suitable class of functions (which will be called test functions) into $\mathbb{R}$. We shall see that this mapping has certain continuity properties. A generalized function is then defined to be a linear mapping on the test functions with these same continuity properties.

Since we intend to use generalized functions to study differential equations, a key question is: how do we define derivatives of such functions? The answer is: by using integration by parts. We shall require $\phi$ to vanish near $\partial\Omega$, so the derivative $\partial f/\partial x_j$ can be defined as the mapping
\begin{equation}\label{eq:distribution-derivative-motivating-pairing}
\phi\mapsto-\int_\Omega f(\bm{x})\dfrac{\partial\phi}{\partial x_j}(\bm{x})\,\mathrm{d}\bm{x}.
\end{equation}
Clearly, this definition requires no differentiability of $f$ in the usual sense; the only differentiability requirement is on $\phi$. We shall therefore choose the test functions to be functions with very nice smoothness properties.
